\section{Implementation}
\label{sec:impl}

Pigs is implemented as a Rust workspace with 14 crates, organized in five layers:
\begin{chinese}
Pigs 以 Rust workspace 形式实现，包含 13 个 crate，组织为五个层次：
\end{chinese}

\begin{center}
\begin{tabular}{lll}
\toprule
Layer & Crates & Role \\
\midrule
0 & pigs-core & Core types: Message, ContentBlock, ApiClient trait \\
1 & pigs-config, pigs-permissions, & Config, permissions, session, prompts \\
  & pigs-session, pigs-prompts & \\
2 & pigs-llm, pigs-tools, pigs-mcp & LLM clients, built-in tools, MCP client \\
3 & pigs-api, pigs-proxy & Phased runtime, HTTP proxy + loopback \\
4 & pigs-cli, pigs-tui, pigs & Agent logic, terminal UI, product binary \\
\bottomrule
\end{tabular}
\end{center}

\subsection{Protocol Codec}
\begin{chinese}
\subsection{协议编解码器}
\end{chinese}

The \texttt{ProtocolCodec} in \texttt{pigs-api} handles parsing, mutation, and response extraction for all three protocols. It validates the ``current user input'' (the last user message, not just the last non-system message) and identifies trailing tool-result groups for continuation matching.
\begin{chinese}
\texttt{pigs-api} 中的 \texttt{ProtocolCodec} 负责所有三种协议的解析、修改和响应提取。它验证"当前用户输入"（最后一条 user 消息，而非仅仅是最后一条非系统消息），并识别尾部的工具结果组以进行续传匹配。
\end{chinese}

Each protocol has dedicated parsing logic:
\begin{chinese}
每种协议具有专用的解析逻辑：
\end{chinese}
\begin{itemize}[nosep]
    \item \textbf{Chat}: messages array with role-based content; tool calls as \texttt{tool\_calls} array; tool results as \texttt{role: "tool"} messages.
    \begin{chinese}
    \textbf{Chat}：具有基于角色内容的 messages 数组；工具调用为 \texttt{tool\_calls} 数组；工具结果为 \texttt{role: "tool"} 消息。
    \end{chinese}
    \item \textbf{Anthropic}: messages with content block arrays; tool calls as \texttt{tool\_use} blocks; tool results as \texttt{tool\_result} blocks within user messages.
    \begin{chinese}
    \textbf{Anthropic}：具有内容块数组的 messages；工具调用为 \texttt{tool\_use} 块；工具结果为 user 消息内的 \texttt{tool\_result} 块。
    \end{chinese}
    \item \textbf{Responses}: input array with typed items; tool calls as \texttt{function\_call} items; tool results as \texttt{function\_call\_output} items.
    \begin{chinese}
    \textbf{Responses}：具有类型化条目的 input 数组；工具调用为 \texttt{function\_call} 条目；工具结果为 \texttt{function\_call\_output} 条目。
    \end{chinese}
\end{itemize}

The codec clones the original body and modifies only the JSON path leading to the current user input, leaving all other fields structurally identical.
\begin{chinese}
编解码器克隆原始 body 并仅修改指向当前用户输入的 JSON 路径，保持所有其他字段在结构上完全一致。
\end{chinese}

\subsection{HTTP Runtime}
\begin{chinese}
\subsection{HTTP 运行时}
\end{chinese}

The \texttt{HttpPhasedRuntime} drives the phase loop. For each phase, it:
\begin{chinese}
\texttt{HttpPhasedRuntime} 驱动阶段循环。对于每个阶段，它：
\end{chinese}
\begin{enumerate}[nosep]
    \item Constructs a phase request via \texttt{phase\_request()} (Pre/Executor/Post prompt + transcript).
    \begin{chinese}
    通过 \texttt{phase\_request()} 构建阶段请求（Pre/Executor/Post 提示 + 记录）。
    \end{chinese}
    \item Sends the request via \texttt{PhaseTransport::send\_streaming()} (streaming) or \texttt{send()} (non-streaming).
    \begin{chinese}
    通过 \texttt{PhaseTransport::send\_streaming()}（流式）或 \texttt{send()}（非流式）发送请求。
    \end{chinese}
    \item Extracts the normalized model output (visible text, tool calls, usage, stop reason).
    \begin{chinese}
    提取归一化的模型输出（可见文本、工具调用、使用量、停止原因）。
    \end{chinese}
    \item Checks for tool calls $\rightarrow$ pauses and stores continuation.
    \begin{chinese}
    检查工具调用 $\rightarrow$ 暂停并存储续传。
    \end{chinese}
    \item Advances the state machine via \texttt{OrchestrationState::advance()}.
    \begin{chinese}
    通过 \texttt{OrchestrationState::advance()} 推进状态机。
    \end{chinese}
\end{enumerate}

The runtime is transport-agnostic: it depends only on the \texttt{PhaseTransport} trait, not on HTTP or axum. This allows the same runtime to be used with any transport implementation.
\begin{chinese}
运行时与传输无关：它仅依赖于 \texttt{PhaseTransport} trait，不依赖 HTTP 或 axum。这使得同一运行时可与任何传输实现配合使用。
\end{chinese}

\subsection{Loopback Transport}
\begin{chinese}
\subsection{回环传输}
\end{chinese}

The \texttt{LoopbackPhaseTransport} in \texttt{pigs-proxy} sends phase subrequests via HTTP to the local proxy. It parses both non-streaming JSON and SSE responses from the upstream provider, aggregating:
\begin{chinese}
\texttt{pigs-proxy} 中的 \texttt{LoopbackPhaseTransport} 通过 HTTP 将阶段子请求发送至本地代理。它解析来自上游提供商的非流式 JSON 和 SSE 响应，聚合：
\end{chinese}
\begin{itemize}[nosep]
    \item Chat: text deltas, tool call deltas (merged by index), finish reason, usage, provider extensions.
    \begin{chinese}
    Chat：文本增量、工具调用增量（按索引合并）、完成原因、使用量、提供商扩展。
    \end{chinese}
    \item Anthropic: content blocks (text, thinking, signature, tool\_use with \texttt{input\_json\_delta} accumulation), stop reason, usage.
    \begin{chinese}
    Anthropic：内容块（text、thinking、signature、tool\_use 及 \texttt{input\_json\_delta} 累积）、停止原因、使用量。
    \end{chinese}
    \item Responses: output items (message, reasoning, function\_call), status, incomplete details, usage.
    \begin{chinese}
    Responses：输出条目（message、reasoning、function\_call）、状态、不完整详情、使用量。
    \end{chinese}
\end{itemize}

Malformed streams (missing termination events, unknown content block indices, open blocks at stream end) are reported as typed errors, not silently accepted.
\begin{chinese}
格式错误的流（缺少终止事件、未知内容块索引、流结束时未关闭的块）被报告为类型化错误，而非被静默接受。
\end{chinese}

\subsection{Body Cleaning}
\begin{chinese}
\subsection{主体清洗}
\end{chinese}

The proxy's \texttt{clean\_empty\_messages} function removes messages with empty content, but preserves assistant messages that contain \texttt{tool\_calls} (even if their content is empty). For messages with empty content and tool calls, the content is replaced with a single space to satisfy upstream providers that reject empty strings.
\begin{chinese}
代理的 \texttt{clean\_empty\_messages} 函数移除内容为空的消息，但保留包含 \texttt{tool\_calls} 的 assistant 消息（即使其内容为空）。对于内容为空且含工具调用的消息，内容被替换为单个空格，以满足上游提供商对空字符串的拒绝要求。
\end{chinese}

\subsection{CLI Integration}
\begin{chinese}
\subsection{CLI 集成}
\end{chinese}

The CLI connects to the local API server via HTTP using \texttt{HttpAgentClient}, which implements the \texttt{ApiClient} trait. The CLI sends requests with the real (non-\texttt{-pig}) model name, so the proxy routes them through passthrough. The CLI's local \texttt{PhasedRuntime} performs Pre$\rightarrow$Executor$\rightarrow$Post orchestration, with tools executed in-process. This design gives the CLI real SSE streaming while maintaining local tool execution.
\begin{chinese}
CLI 通过 HTTP 使用 \texttt{HttpAgentClient} 连接本地 API 服务器，后者实现了 \texttt{ApiClient} trait。CLI 以真实（非 \texttt{-pig}）模型名发送请求，因此代理通过透传路径路由。CLI 本地的 \texttt{PhasedRuntime} 执行 Pre$\rightarrow$Executor$\rightarrow$Post 编排，工具在进程内执行。此设计使 CLI 在保持本地工具执行的同时获得真实的 SSE 流式传输。
\end{chinese}
